% =====================================================================
%  CHƯƠNG 3: PHÂN TÍCH VÀ THIẾT KẾ HỆ THỐNG
% =====================================================================
\chapter{Phân tích và thiết kế hệ thống}
\label{chap:design}

\section{Kiến trúc tổng thể}

Hệ thống được thiết kế thành một pipeline gồm bốn giai đoạn nối tiếp, từ dữ liệu thô
tới danh sách gợi ý (Hình~\ref{fig:pipeline}): (1) tiền xử lý dữ liệu lớn bằng Spark
và lưu Parquet; (2) mã hóa "trải nghiệm" người dùng và tổng hợp lịch sử có trọng số
thời gian; (3) huấn luyện mô hình hai tháp với học tương phản; (4) đánh giá truy hồi
trên toàn bộ kho sản phẩm.

Hình~\ref{fig:pipeline} là sơ đồ khối chi tiết toàn bộ luồng -- từ dữ liệu thô, qua
xử lý dữ liệu lớn, trích xuất đặc trưng, tới mô hình và đánh giá -- kèm các tệp/đối
tượng trung gian (Parquet, vector lịch sử) chuyển giao giữa các giai đoạn.

\begin{figure}[H]
\centering
\begin{tikzpicture}[
    font=\scriptsize,
    io/.style={rectangle, rounded corners=2pt, draw=gray!55, fill=gray!12,
        line width=0.6pt, text width=2.0cm, minimum height=0.65cm, align=center},
    proc/.style={rectangle, rounded corners=2pt, draw=blue!55, fill=blue!12,
        line width=0.7pt, text width=6.6cm, minimum height=0.62cm, align=center},
    art/.style={rectangle, rounded corners=2pt, draw=violet!60, fill=violet!8,
        line width=0.7pt, text width=6.6cm, minimum height=0.55cm, align=center},
    mdl/.style={rectangle, rounded corners=2pt, draw=teal!60, fill=teal!10,
        line width=0.7pt, text width=3.0cm, minimum height=0.78cm, align=center},
    a/.style={-{Stealth[length=1.8mm]}, line width=0.7pt, gray!65},
    plab/.style={font=\scriptsize\bfseries, anchor=west}
]
% ===== GIAI ĐOẠN 1: XỬ LÝ DỮ LIỆU =====
\node[io] (rev)   at (-2.6,0) {Reviews\\(JSONL)};
\node[io] (meta)  at (0,0)    {Metadata\\(JSONL)};
\node[io] (split) at (2.6,0)  {Split CSV\\train/val/test};
\node[proc] (spark) at (0,-1.25) {\textbf{Apache Spark:} khử trùng lặp $\cdot$ broadcast join $\cdot$ dựng \texttt{semantic\_text} $\cdot$ \texttt{log1p}};
\node[art]  (parq)  at (0,-2.35) {\textbf{Parquet:} train / val / test \;+\; item};
\draw[a] (rev)--(spark); \draw[a] (meta)--(spark); \draw[a] (split)--(spark);
\draw[a] (spark)--(parq);
% ===== GIAI ĐOẠN 2: TRÍCH XUẤT ĐẶC TRƯNG =====
\node[proc, draw=orange!65, fill=orange!12] (enc) at (0,-4.25)
    {\textbf{Mã hóa trải nghiệm} bằng MiniLM (đóng băng) + mean pooling};
\node[proc, draw=orange!65, fill=orange!12] (agg) at (0,-5.35)
    {\textbf{Tổng hợp trọng số thời gian} (recency $\times$ rating)};
\node[art] (uvec) at (0,-6.45) {\textbf{Vector lịch sử người dùng} (384-d)};
\draw[a] (parq)--(enc); \draw[a] (enc)--(agg); \draw[a] (agg)--(uvec);
% ===== GIAI ĐOẠN 3: MÔ HÌNH =====
\node[mdl] (ut) at (-2.4,-8.7) {Tháp User\\(MLP $384{\to}512{\to}384$)};
\node[mdl, draw=gray!55, fill=gray!12] (it) at (2.4,-8.7) {Tháp Item\\(MiniLM, đóng băng)};
\node[mdl, draw=orange!70, fill=orange!14] (loss) at (0,-10.1) {Cosine + InfoNCE\\(in-batch negatives)};
\draw[a] (uvec)--(ut);
\draw[a, gray!55] (parq.east) to[out=0,in=90] (it.north);  % van ban item tu Parquet item
\draw[a] (ut)--(loss); \draw[a, gray!55] (it)--(loss);
\draw[a, dashed, orange!70] (loss.west) to[out=180,in=-90] node[left, font=\tiny] {cập nhật} (ut.south);
% ===== GIAI ĐOẠN 4: ĐÁNH GIÁ =====
\node[mdl, draw=teal!60, fill=teal!12] (metric) at (0,-12.3)
    {\textbf{Đánh giá full-catalog} ($\sim$17K item): Recall@K, NDCG@K};
\draw[a, teal!60] (loss)--(metric);
% ===== BĂNG NỀN CÁC GIAI ĐOẠN =====
\begin{scope}[on background layer]
\node[draw=blue!50, rounded corners=4pt, fill=blue!4, line width=0.8pt,
      fit=(rev)(split)(parq), inner sep=6pt] (b1) {};
\node[draw=orange!55, rounded corners=4pt, fill=orange!4, line width=0.8pt,
      fit=(enc)(uvec), inner sep=6pt] (b2) {};
\node[draw=teal!55, rounded corners=4pt, fill=teal!4, line width=0.8pt,
      fit=(ut)(it)(loss), inner sep=6pt] (b3) {};
\node[draw=teal!60, rounded corners=4pt, fill=teal!7, line width=0.8pt,
      fit=(metric), inner sep=5pt] (b4) {};
\end{scope}
\node[plab, blue!55!black]   at ([yshift=8pt]b1.north west) {1. XỬ LÝ DỮ LIỆU (Spark)};
\node[plab, orange!60!black] at ([yshift=8pt]b2.north west) {2. TRÍCH XUẤT ĐẶC TRƯNG};
\node[plab, teal!55!black]   at ([yshift=8pt]b3.north west) {3. MÔ HÌNH TWO-TOWER};
\node[plab, teal!55!black]   at ([yshift=7pt]b4.north west) {4. ĐÁNH GIÁ};
\end{tikzpicture}
\caption{Sơ đồ khối end-to-end: từ xử lý dữ liệu đến mô hình và đánh giá}
\label{fig:pipeline}
\end{figure}

\section{Dữ liệu và tiền xử lý bằng PySpark}

Dữ liệu sử dụng là \textbf{Amazon Reviews 2023} (nhóm McAuley Lab, UCSD), danh mục
\textit{Software}, gồm hai phần: tập \textbf{tương tác} (review của người dùng, các
trường \texttt{user\_id}, \texttt{parent\_asin}, \texttt{rating}, \texttt{title},
\texttt{text}, \texttt{timestamp}) và tập \textbf{metadata sản phẩm} (tiêu đề, tính
năng, mô tả, danh mục, giá, điểm trung bình\ldots). Quy mô bài toán được tổng hợp
trong Bảng~\ref{tab:stats}.

\begin{table}[H]
\centering
\caption{Thống kê quy mô bộ dữ liệu Amazon Software}
\label{tab:stats}
\renewcommand{\arraystretch}{1.3}
\begin{tabularx}{0.86\textwidth}{@{}X r@{}}
\toprule
\textbf{Thuộc tính} & \textbf{Giá trị} \\
\midrule
Số người dùng                         & $146.396$ \\
Số sản phẩm phần mềm độc nhất          & $> 17.000$ \\
Số tương tác huấn luyện (train)        & $\approx 984.000$ \\
Số mẫu kiểm định (validation)          & $146.396$ (1 / người dùng) \\
Số mẫu kiểm thử (test)                 & $146.396$ (1 / người dùng) \\
Số chiều vector biểu diễn              & $384$ \\
Thang điểm đánh giá                    & $1$ -- $5$ sao \\
\bottomrule
\end{tabularx}
\end{table}

Theo giao thức leave-last-out, mỗi người dùng đóng góp đúng một mẫu cho tập kiểm
định và một mẫu cho tập kiểm thử, phần còn lại thuộc tập huấn luyện. Phân phối số
tương tác trên mỗi người dùng (Hình~\ref{fig:userint}) lệch phải mạnh -- đa số người
dùng chỉ có vài tương tác, một thiểu số rất tích cực -- phản ánh tính thưa thớt và
đuôi dài đặc trưng của dữ liệu thương mại điện tử.

\begin{figure}[H]
\centering
\includegraphics[width=0.72\textwidth]{user_interactions.png}
\caption{Phân phối số tương tác trên mỗi người dùng}
\label{fig:userint}
\end{figure}

\subsection{Làm sạch và kiểm tra chất lượng dữ liệu}

Dữ liệu thô từ Amazon chứa nhiều giá trị thiếu và bản ghi trùng. Quy trình làm sạch
trên Spark gồm: (i) \textbf{khai báo lược đồ tường minh} (\texttt{StructType}) khi
đọc metadata để kiểm soát kiểu dữ liệu của các trường phức tạp (mảng, map); (ii)
\textbf{đếm giá trị thiếu} từng cột để đánh giá chất lượng; (iii) \textbf{khử trùng
lặp} bằng \texttt{dropDuplicates}; và (iv) \textbf{ghép nối} tập chia chuẩn với dữ
liệu đầy đủ bằng \texttt{broadcast join} -- kỹ thuật phát tán bảng nhỏ tới mọi
executor, tránh thao tác shuffle tốn kém trên bảng lớn.

\begin{lstlisting}[language=Python, caption=Kiểm tra null và ghép nối broadcast trên Spark]
from pyspark.sql.functions import col, count, when, isnan, broadcast

# Dem so gia tri thieu tren tung cot
df.select([
    count(when(col(c).isNull(), c)).alias(c) for c in df.columns
]).show()

# Khu trung lap
df = df.dropDuplicates()

# Ghep tap split chuan (nho) voi du lieu day du (lon) bang broadcast join
df_train_ok = broadcast(df_train).join(
    df, on=["user_id", "parent_asin", "timestamp"], how="left")
\end{lstlisting}

Các trường số như \texttt{price} và \texttt{rating\_number} có phân phối lệch phải
rất mạnh nên được biến đổi \texttt{log1p} để nén khoảng giá trị, giúp ổn định khi
đưa vào mô hình. Cuối cùng, tập tương tác được lọc để chỉ giữ lại các sản phẩm tồn
tại trong bảng metadata, bảo đảm mọi sản phẩm đều có văn bản ngữ nghĩa.

\subsection{Dựng văn bản ngữ nghĩa cho sản phẩm}

Bước quan trọng nhất của tiền xử lý là tổng hợp các trường metadata rời rạc thành một
chuỗi văn bản ngữ nghĩa thống nhất (\texttt{semantic\_text}) để làm đầu vào cho mô
hình ngôn ngữ. Các trường dạng mảng được ghép bằng \texttt{concat\_ws}, trường
\texttt{details} dạng map được duyệt thành các dòng "khóa: giá trị":

\begin{lstlisting}[language=Python, caption=Dựng semantic\_text từ metadata bằng PySpark]
from pyspark.sql import functions as F

# Ghep cac truong mang (features, description, categories) thanh text
for c in ["features", "description", "categories"]:
    df_i = df_i.withColumn(f"{c}_text",
        F.when(F.col(c).isNotNull(), F.concat_ws(" ", F.col(c))).otherwise(F.lit("")))

# Tong hop semantic_text co cau truc theo tung khoi
df_i = df_i.withColumn("semantic_text", F.concat_ws("\n\n",
    F.when(F.col("title") != "",         F.concat(F.lit("TITLE\n"), F.col("title"))),
    F.when(F.col("features_text") != "", F.concat(F.lit("FEATURES\n"), F.col("features_text"))),
    F.when(F.col("description_text") != "", F.concat(F.lit("DESCRIPTION\n"), F.col("description_text"))),
    F.when(F.col("categories_text") != "", F.concat(F.lit("CATEGORIES\n"), F.col("categories_text")))))

# Chuan hoa truong so: log1p cho gia va so luong danh gia (giam lech phai)
df_i = df_i.withColumn("price", F.log1p(F.col("price")))
df_i = df_i.withColumn("rating_number", F.log1p(F.col("rating_number")))
\end{lstlisting}

Một ví dụ \texttt{semantic\_text} sau khi tổng hợp có dạng như sau:

\begin{lstlisting}[caption=Ví dụ văn bản ngữ nghĩa của một sản phẩm phần mềm]
TITLE
Norton 360 Deluxe 2024, Antivirus software for 5 Devices

FEATURES
Real-time protection against malware, ransomware and phishing
Secure VPN for online privacy; 50 GB cloud backup

CATEGORIES
Software | Antivirus & Security | Internet Security Suites
\end{lstlisting}

Khi đưa vào mô hình ngôn ngữ, toàn bộ khối văn bản này được nén thành một vector
$384$ chiều duy nhất, "hiểu" được rằng đây là phần mềm bảo mật -- nhờ đó các sản
phẩm cùng loại sẽ nằm gần nhau trong không gian biểu diễn.

Dữ liệu sau xử lý được lưu thành các tệp Parquet riêng cho tập huấn luyện, kiểm định,
kiểm thử và bảng sản phẩm. Phân phối điểm đánh giá (Hình~\ref{fig:rating}) lệch mạnh
về phía điểm cao -- đặc trưng điển hình của dữ liệu thương mại điện tử. Độ dài
\texttt{semantic\_text} (Hình~\ref{fig:textlen}) có trung vị khoảng $70$--$80$ token;
mô hình giới hạn \texttt{max\_length = 128} để cân bằng giữa thông tin và tốc độ.

\begin{figure}[H]
\centering
\begin{subfigure}{0.49\textwidth}
    \centering\includegraphics[width=\textwidth]{rating_distribution.png}
    \caption{Phân phối điểm đánh giá}
    \label{fig:rating}
\end{subfigure}
\hfill
\begin{subfigure}{0.49\textwidth}
    \centering\includegraphics[width=\textwidth]{text_length.png}
    \caption{Độ dài semantic\_text}
    \label{fig:textlen}
\end{subfigure}
\caption{Một số đặc trưng của dữ liệu Amazon Software}
\label{fig:data-eda}
\end{figure}

\section{Chiến lược chia dữ liệu theo thời gian}

Để mô phỏng sát kịch bản thực tế (dự đoán hành vi \textit{tương lai} từ
\textit{quá khứ}), đồ án áp dụng phương pháp \textbf{cắt chuỗi hành vi theo thời
gian} (leave-last-out) cho mỗi người dùng:

\begin{itemize}[leftmargin=2em, itemsep=2pt]
    \item \textbf{Tập kiểm thử (test):} sản phẩm được tương tác \textit{cuối cùng}.
    \item \textbf{Tập kiểm định (validation):} sản phẩm \textit{kề cuối}.
    \item \textbf{Tập huấn luyện (train):} toàn bộ các tương tác \textit{trước đó}.
\end{itemize}

Cách chia này nghiêm ngặt hơn chia ngẫu nhiên vì tuyệt đối không để lộ thông tin
tương lai vào quá khứ (Hình~\ref{fig:loo}). Lịch sử của tập kiểm định kế thừa từ tập
huấn luyện; lịch sử của tập kiểm thử được tính trên hợp của tập huấn luyện và kiểm
định.

\begin{figure}[H]
\centering
\begin{tikzpicture}[font=\footnotesize,
    tr/.style={rectangle, rounded corners=2pt, draw=blue!55, fill=blue!10, minimum width=1.15cm, minimum height=0.8cm},
    va/.style={rectangle, rounded corners=2pt, draw=orange!70, fill=orange!15, minimum width=1.15cm, minimum height=0.8cm},
    te/.style={rectangle, rounded corners=2pt, draw=red!65, fill=red!12, minimum width=1.15cm, minimum height=0.8cm}]
\node[tr] (t1) {$i_1$};
\node[tr, right=0.18cm of t1] (t2) {$i_2$};
\node[right=0.18cm of t2] (dots) {$\cdots$};
\node[tr, right=0.18cm of dots] (t3) {$i_{n-2}$};
\node[va, right=0.18cm of t3] (t4) {$i_{n-1}$};
\node[te, right=0.18cm of t4] (t5) {$i_{n}$};
\draw[-{Stealth[length=2.4mm]}, line width=0.8pt, gray!70]
    ([yshift=-0.65cm]t1.south west) -- ([yshift=-0.65cm]t5.south east)
    node[midway, below, font=\scriptsize] {Thời gian};
\node[below=0.02cm of t1, font=\scriptsize, blue!55!black] {};
\draw[decorate, decoration={brace, amplitude=4pt}, blue!55]
    (t1.north west) -- (t3.north east) node[midway, above=3pt, font=\scriptsize\bfseries, blue!55!black] {Train};
\node[above=0.5cm of t4, font=\scriptsize\bfseries, orange!70!black] {Val};
\node[above=0.5cm of t5, font=\scriptsize\bfseries, red!65!black] {Test};
\draw[-{Stealth[length=1.6mm]}, orange!70] ([yshift=0.32cm]t4.north) -- (t4.north);
\draw[-{Stealth[length=1.6mm]}, red!65] ([yshift=0.32cm]t5.north) -- (t5.north);
\end{tikzpicture}
\caption{Chia dữ liệu leave-last-out: mục cuối làm test, kề cuối làm val, còn lại train}
\label{fig:loo}
\end{figure}

\section{Mã hóa trải nghiệm và tổng hợp lịch sử người dùng}

\subsection{Mã hóa trải nghiệm}

Trước hết, mỗi sản phẩm được biểu diễn bằng một chuỗi văn bản tổng hợp
(\texttt{item\_text\_map}) ghép \texttt{semantic\_text} với các thuộc tính bổ sung
(danh mục, cửa hàng, giá, điểm và số lượng đánh giá):

\begin{lstlisting}[language=Python, caption=Dựng từ điển văn bản sản phẩm item\_text\_map]
item_text_map = {}
for r in item_df.itertuples():
    item_text_map[r.parent_asin] = (
        f"{safe(r.semantic_text)}\n"
        f"Category: {safe(r.main_category)}\nStore: {safe(r.store)}\n"
        f"Price: {safe(r.price)}\nAverage Rating: {safe(r.average_rating)}\n"
        f"Rating Count: {safe(r.rating_number)}")
\end{lstlisting}

Mỗi tương tác của người dùng được biến thành một "trải nghiệm" bằng cách ghép văn bản
sản phẩm trên với tiêu đề và nội dung review, rồi mã hóa qua mô hình MiniLM (đóng băng)
theo từng batch trên GPU:

\begin{lstlisting}[language=Python, caption=Mã hóa trải nghiệm người dùng hàng loạt trên GPU]
def precompute_user_experience_vectors(df, batch_size=256):
    base = df['parent_asin'].map(item_text_map).fillna('None')
    combined = (base + "\nUser Review Title: " + df['title'].fillna('None')
                     + "\nUser Review Content: " + df['text'].fillna('None')).tolist()
    vectors = []
    with torch.no_grad():
        for i in range(0, len(combined), batch_size):
            enc = tokenizer(combined[i:i+batch_size], padding=True, truncation=True,
                            max_length=128, return_tensors="pt").to(device)
            out = base_model(**enc)
            emb = F.normalize(mean_pooling(out, enc["attention_mask"]), dim=1)
            vectors.append(emb.cpu().numpy())
    return np.vstack(vectors)
\end{lstlisting}

\subsection{Tổng hợp lịch sử với trọng số thời gian}

Vector đại diện cho một người dùng là \textit{trung bình có trọng số} của các vector
trải nghiệm trong lịch sử. Trọng số kết hợp hai yếu tố: (i) \textbf{độ gần đây}
(recency) -- tương tác càng mới càng quan trọng, theo hàm suy giảm mũ
$e^{-(n-i)/n}$; và (ii) \textbf{mức độ ưa thích} -- điểm rating càng cao trọng số
càng lớn, theo hệ số $(1 + \text{rating}/5)$:

\begin{equation}
    \mathbf{u} = \frac{\sum_{i=1}^{n} w_i\, \mathbf{e}_i}{\sum_{i=1}^{n} w_i},
    \qquad
    w_i = e^{-(n-i)/n} \cdot \left(1 + \frac{\text{rating}_i}{5}\right),
\end{equation}

trong đó $\mathbf{e}_i$ là vector trải nghiệm thứ $i$ và $n$ là số tương tác của
người dùng. Vector kết quả được chuẩn hóa $L_2$. Hình~\ref{fig:timedecay} minh họa
cách trọng số tăng dần về phía các tương tác gần đây nhất. Toàn bộ phép tổng hợp này
chạy trên CPU/RAM (gom nhóm theo \texttt{user\_id}), độc lập với GPU, nên rất nhẹ và
nhanh.

\begin{lstlisting}[language=Python, caption=Tổng hợp vector lịch sử có trọng số thời gian]
def aggregate_user_history(df, experience_vectors):
    df['_vec'] = list(experience_vectors)
    user_hist = {}
    for uid, g in df.groupby("user_id"):
        g = g.reset_index(drop=True); n = len(g)
        vectors = np.array(g['_vec'].tolist())
        idx = np.arange(n)
        recency = np.exp(-(n - idx) / max(n, 1))          # cang moi cang nang
        ratings = g['rating'].fillna(0).to_numpy()
        weights = (recency * (1 + ratings / 5.0)).reshape(-1, 1)
        u = np.sum(vectors * weights, axis=0) / np.sum(weights)
        user_hist[uid] = u / (np.linalg.norm(u) + 1e-9)   # chuan hoa L2
    return user_hist
\end{lstlisting}

Phép \textbf{mean pooling có mặt nạ} (dùng ở cả khâu mã hóa) lấy trung bình các vector
token, bỏ qua token đệm nhờ \texttt{attention\_mask}:

\begin{lstlisting}[language=Python, caption=Mean pooling có mặt nạ]
def mean_pooling(out, mask):
    token = out.last_hidden_state                      # [B, L, 384]
    mask = mask.unsqueeze(-1).expand(token.size()).float()
    return (token * mask).sum(1) / torch.clamp(mask.sum(1), min=1e-9)
\end{lstlisting}

\begin{figure}[H]
\centering
\includegraphics[width=0.74\textwidth]{time_decay.png}
\caption{Trọng số thời gian khi tổng hợp vector lịch sử người dùng ($n = 15$)}
\label{fig:timedecay}
\end{figure}

Toàn bộ luồng tổng hợp được tóm tắt trong Hình~\ref{fig:aggflow}: các vector trải
nghiệm được nhân với trọng số tương ứng, cộng lại rồi chuẩn hóa thành một vector
người dùng duy nhất.

\begin{figure}[H]
\centering
\begin{tikzpicture}[font=\footnotesize, node distance=0.7cm,
    e/.style={rectangle, rounded corners=2pt, draw=blue!55, fill=blue!9, minimum width=1.3cm, minimum height=0.7cm},
    op/.style={circle, draw=teal!60, fill=teal!12, minimum size=0.8cm},
    o/.style={rectangle, rounded corners=2pt, draw=orange!70, fill=orange!12, minimum width=1.8cm, minimum height=0.7cm},
    a/.style={-{Stealth[length=2mm]}, line width=0.7pt, blue!55}]
\node[e] (e1) {$\mathbf{e}_1$};
\node[e, below=0.3cm of e1] (e2) {$\mathbf{e}_2$};
\node[below=0.18cm of e2, font=\scriptsize] (ed) {$\vdots$};
\node[e, below=0.25cm of ed] (en) {$\mathbf{e}_n$};
\node[op, right=1.6cm of e2] (sum) {$\sum$};
\node[o, right=1.3cm of sum] (norm) {L2 Normalize};
\node[o, right=0.9cm of norm, fill=orange!18] (u) {$\mathbf{u}$};
\draw[a] (e1) -- node[above, font=\scriptsize, pos=0.45] {$\times w_1$} (sum);
\draw[a] (e2) -- node[above, font=\scriptsize, pos=0.5] {$\times w_2$} (sum);
\draw[a] (en) -- node[below, font=\scriptsize, pos=0.45] {$\times w_n$} (sum);
\draw[a] (sum) -- (norm);
\draw[a] (norm) -- (u);
\node[below=0.55cm of sum, font=\scriptsize, gray!55!black, text width=3cm, align=center]
    {$w_i = e^{-(n-i)/n}(1+\text{rating}_i/5)$};
\end{tikzpicture}
\caption{Luồng tổng hợp vector lịch sử người dùng từ các vector trải nghiệm}
\label{fig:aggflow}
\end{figure}

\section{Thiết kế mô hình hai tháp}

Mô hình \texttt{TwoTowerVectorModel} hiện thực hóa kiến trúc ở
Mục~\ref{sec:twotower}. \textbf{Tháp người dùng} là một mạng MLP biến đổi vector
lịch sử $384$ chiều, có \texttt{BatchNorm} để ổn định gradient và \texttt{Dropout}
để chống quá khớp. \textbf{Tháp sản phẩm} là MiniLM \textbf{đóng băng}, chỉ làm
nhiệm vụ mã hóa văn bản.

\begin{lstlisting}[language=Python, caption=Mô hình hai tháp (tháp sản phẩm đóng băng)]
class TwoTowerVectorModel(nn.Module):
    def __init__(self, model_name):
        super().__init__()
        # Thap User: MLP co the huan luyen
        self.user_projection = nn.Sequential(
            nn.Linear(384, 512), nn.BatchNorm1d(512), nn.ReLU(),
            nn.Dropout(0.5), nn.Linear(512, 384))
        # Thap Item: dong bang hoan toan
        self.item_encoder = AutoModel.from_pretrained(model_name)
        for p in self.item_encoder.parameters():
            p.requires_grad = False

    def forward_user(self, user_vectors):
        return F.normalize(self.user_projection(user_vectors), dim=1)

    def forward_item(self, tokens):
        with torch.no_grad():
            out = self.item_encoder(**tokens)
            emb = self.mean_pooling(out, tokens["attention_mask"])
        return F.normalize(emb, dim=1)
\end{lstlisting}

\subsection{Sơ đồ kiến trúc chi tiết}

Hình~\ref{fig:model-detail} mô tả chi tiết từng tầng của mô hình cùng kích thước
tensor tương ứng ($B$ là kích thước batch, $L$ là độ dài chuỗi token). Tháp người
dùng biến đổi vector lịch sử qua một khối MLP, còn tháp sản phẩm đưa chuỗi token qua
$6$ tầng Transformer của MiniLM rồi gộp trung bình. Hai vector đầu ra cùng kích
thước $384$ chiều được chuẩn hóa và đem so khớp bằng tích vô hướng.

\begin{figure}[H]
\centering
\begin{tikzpicture}[
    font=\scriptsize,
    layer/.style={rectangle, rounded corners=2pt, draw=blue!55, fill=blue!8,
        line width=0.7pt, text width=2.75cm, minimum height=0.6cm, align=center},
    frozen/.style={rectangle, rounded corners=2pt, draw=gray!60, fill=gray!12,
        line width=0.7pt, text width=2.75cm, minimum height=0.6cm, align=center},
    io/.style={rectangle, rounded corners=2pt, draw=teal!60, fill=teal!10,
        line width=0.7pt, text width=2.75cm, minimum height=0.6cm, align=center},
    sh/.style={font=\scriptsize\itshape, text=gray!55!black},
    a/.style={-{Stealth[length=1.6mm]}, line width=0.6pt, blue!55},
    ag/.style={-{Stealth[length=1.6mm]}, line width=0.6pt, gray!60}
]
% ===== USER TOWER (left) =====
\node[io] (uin) {Vector lịch sử người dùng};
\node[sh, right=0.15cm of uin] {[B, 384]};
\node[layer, above=0.4cm of uin] (l1) {Linear $384 \to 512$};
\node[sh, right=0.15cm of l1] {[B, 512]};
\node[layer, above=0.35cm of l1] (l2) {BatchNorm1d(512)};
\node[layer, above=0.35cm of l2] (l3) {ReLU};
\node[layer, above=0.35cm of l3] (l4) {Dropout($0.5$)};
\node[layer, above=0.35cm of l4] (l5) {Linear $512 \to 384$};
\node[sh, right=0.15cm of l5] {[B, 384]};
\node[io, above=0.4cm of l5] (un) {L2 Normalize};
\draw[a] (uin)--(l1); \draw[a] (l1)--(l2); \draw[a] (l2)--(l3);
\draw[a] (l3)--(l4); \draw[a] (l4)--(l5); \draw[a] (l5)--(un);

% ===== ITEM TOWER (right) =====
\node[io, right=3.0cm of uin] (iin) {Tokens semantic\_text};
\node[sh, right=0.15cm of iin] {[B, L]};
\node[frozen, above=0.4cm of iin] (e1) {Embedding};
\node[sh, right=0.15cm of e1] {[B, L, 384]};
\node[frozen, above=0.35cm of e1] (e2) {Transformer $\times\,6$\\(self-attention + FFN)};
\node[sh, right=0.15cm of e2] {[B, L, 384]};
\node[frozen, above=0.62cm of e2] (e3) {Mean Pooling (có mặt nạ)};
\node[sh, right=0.15cm of e3] {[B, 384]};
\node[io, above=0.4cm of e3] (inr) {L2 Normalize};
\draw[ag] (iin)--(e1); \draw[ag] (e1)--(e2); \draw[ag] (e2)--(e3); \draw[ag] (e3)--(inr);
\node[font=\scriptsize\bfseries, below=0.12cm of iin, text=gray!55!black]
    {Tháp Item (đóng băng)};
\node[font=\scriptsize\bfseries, below=0.12cm of uin, text=blue!55!black]
    {Tháp User (huấn luyện)};

% ===== TOP: similarity =====
\node[rectangle, rounded corners=3pt, draw=orange!70, fill=orange!12,
      line width=0.9pt, text width=4.8cm, align=center, minimum height=0.75cm,
      above=1.15cm of $(un)!0.5!(inr)$] (sim)
      {Ma trận tương đồng $\mathbf{U}\mathbf{V}^\top$ $\;$[B, B]\\$\to$ InfoNCE Loss};
\draw[a] (un) |- (sim);
\draw[ag] (inr) |- (sim);
\end{tikzpicture}
\caption{Sơ đồ kiến trúc chi tiết của mô hình hai tháp kèm kích thước tensor}
\label{fig:model-detail}
\end{figure}

\subsection{Phân tích số tham số}

Một ưu điểm nổi bật của thiết kế là \textbf{hiệu quả tham số}. Bảng~\ref{tab:params}
liệt kê số tham số của tháp người dùng -- chỉ khoảng $395.000$ tham số khả huấn
luyện, trong khi tháp sản phẩm (MiniLM) có $\sim 22{,}7$ triệu tham số nhưng được
\textbf{đóng băng hoàn toàn}. Như vậy, chỉ khoảng \textbf{$1{,}7\%$} tổng số tham số
thực sự được cập nhật (Hình~\ref{fig:params}), giúp giảm mạnh chi phí bộ nhớ và thời
gian huấn luyện.

\begin{table}[H]
\centering
\caption{Số tham số khả huấn luyện của tháp người dùng}
\label{tab:params}
\renewcommand{\arraystretch}{1.25}
\begin{tabularx}{0.82\textwidth}{@{}X r@{}}
\toprule
\textbf{Tầng} & \textbf{Số tham số} \\
\midrule
Linear $384 \to 512$        & $197.120$ \\
BatchNorm1d$(512)$          & $1.024$ \\
Linear $512 \to 384$        & $196.992$ \\
\midrule
\textbf{Tổng (khả huấn luyện)} & $\mathbf{395.136}$ \\
Tháp Item (MiniLM, đóng băng)  & $\approx 22.700.000$ \\
\bottomrule
\end{tabularx}
\end{table}

\begin{figure}[H]
\centering
\includegraphics[width=0.62\textwidth]{param_breakdown.png}
\caption{Phân bổ tham số giữa hai tháp (chỉ $\sim 1{,}7\%$ được huấn luyện)}
\label{fig:params}
\end{figure}

\section{Hàm mất mát và cấu hình huấn luyện}

Mô hình được huấn luyện bằng hàm mất mát \textbf{InfoNCE} (Mục~\ref{sec:infonce})
với mẫu âm trong batch. Ma trận tương đồng giữa mọi người dùng và mọi sản phẩm trong
batch được tính, nhãn đúng nằm trên đường chéo:

\begin{lstlisting}[language=Python, caption=Hàm mất mát InfoNCE với mẫu âm trong batch]
class InfoNCE(nn.Module):
    def __init__(self, temperature=0.1):
        super().__init__()
        self.temperature = temperature
        self.ce = nn.CrossEntropyLoss()

    def forward(self, user_out, item_out):
        sim = torch.matmul(user_out, item_out.T) / self.temperature
        labels = torch.arange(user_out.size(0), device=user_out.device)
        return self.ce(sim, labels)

loss_fn = InfoNCE(temperature=0.1)
optimizer = torch.optim.AdamW(model.parameters(), lr=1e-4, weight_decay=5e-2)
\end{lstlisting}

Mô hình được tối ưu bằng \texttt{AdamW} (learning rate $10^{-4}$, weight decay
$0{,}05$). Quá trình huấn luyện áp dụng \textbf{tính toán độ chính xác hỗn hợp}
(mixed precision) để tiết kiệm bộ nhớ GPU và tăng tốc, kèm cơ chế dừng sớm (early
stopping) dựa trên mất mát kiểm định. Chi tiết cấu hình và kết quả được trình bày ở
Chương~\ref{chap:exp}.
